\chapter{Contexto y punto de partida}
El detector T0 de SHiP se estudia aquí mediante una barra de centelleador plástico con lectura en los extremos y en la cara superior. La geometría, el mapeo de sensores y los identificadores globales están definidos en el esquema de EXEC\_46. Los antecedentes experimentales principales son Betancourt et al.~\cite{Betancourt2020} y Blondel et al.~\cite{Blondel2018}.

La campaña inicial registró una no-linealidad de $T_0(x)$ y un remanente tabulado en Step~5. El inventario de materiales, posiciones y eventos se conserva en los CSV de Step~1 y Step~2. Los cuatro mecanismos candidatos de la auditoría inicial se mantuvieron como hipótesis de trabajo; este reporte tabula las pruebas realizadas sin convertirlas en conclusiones adicionales.

La campaña se compone de tres materiales y siete posiciones. La campaña antigua usa el modelo de índice constante; la campaña v2 usa los runtimes corregidos para EJ--200 y EJ--204, con EJ--230 como control constante.
